%==============================================================================
% Chapitre 18 - Alimentation et cycle de fonctionnement
%==============================================================================

\section{Introduction}

Le système d'alimentation assure l'approvisionnement de l'arme en munitions.

Son rôle consiste à présenter chaque cartouche dans une position permettant
son introduction dans la chambre puis la poursuite du cycle de tir.

La qualité de ce système influence directement :

\begin{itemize}
    \item la fiabilité ;
    \item la cadence de tir ;
    \item l'autonomie ;
    \item la sécurité de fonctionnement.
\end{itemize}

\begin{aretenir}

Une arme extrêmement précise devient inutile si elle ne peut pas alimenter
correctement les munitions.

\end{aretenir}

\section{Fonction d'alimentation}

Le système d'alimentation doit assurer plusieurs fonctions :

\begin{itemize}
    \item stocker les munitions ;
    \item les présenter correctement ;
    \item les acheminer vers la chambre ;
    \item maintenir une alimentation régulière.
\end{itemize}

Ces fonctions doivent être réalisées malgré :

\begin{itemize}
    \item les vibrations ;
    \item les chocs ;
    \item l'encrassement ;
    \item les variations de position de l'arme.
\end{itemize}

\section{Chargeurs}

Le chargeur est aujourd'hui la solution la plus répandue.

Il contient généralement :

\begin{itemize}
    \item un corps ;
    \item un ressort ;
    \item une planchette élévatrice ;
    \item des lèvres de retenue.
\end{itemize}

Le ressort pousse les cartouches vers le haut afin de permettre leur
présentation.

\subsection{Avantages}

\begin{itemize}
    \item simplicité ;
    \item faible masse ;
    \item rechargement rapide ;
    \item fabrication économique.
\end{itemize}

\subsection{Inconvénients}

\begin{itemize}
    \item capacité limitée ;
    \item sensibilité aux déformations ;
    \item sensibilité à certains encrassements.
\end{itemize}

\section{Bandes de munitions}

Les mitrailleuses utilisent fréquemment des bandes de munitions.

Chaque cartouche est maintenue individuellement dans un maillon.

Cette solution permet :

\begin{itemize}
    \item une très grande capacité ;
    \item une alimentation continue ;
    \item une bonne adaptation aux armes automatiques.
\end{itemize}

\begin{exemple}

La MAG58 utilise une alimentation par bande désintégrable.

Chaque maillon est séparé du reste de la bande lors du cycle
d'alimentation.

\end{exemple}

\section{Tambours et magasins spéciaux}

Certaines armes utilisent :

\begin{itemize}
    \item des tambours ;
    \item des magasins circulaires ;
    \item des magasins hélicoïdaux ;
    \item des systèmes spécifiques.
\end{itemize}

Ces solutions visent généralement à augmenter la capacité disponible.

\section{Cycle de fonctionnement}

Le cycle de fonctionnement décrit l'ensemble des opérations réalisées entre
deux tirs successifs.

Il comprend généralement :

\begin{enumerate}
    \item alimentation ;
    \item chambrage ;
    \item verrouillage ;
    \item percussion ;
    \item départ du coup ;
    \item déverrouillage ;
    \item extraction ;
    \item éjection ;
    \item réarmement.
\end{enumerate}

\section{Chambrage}

Le chambrage consiste à introduire une cartouche dans la chambre.

Cette opération doit être réalisée avec précision afin d'éviter :

\begin{itemize}
    \item les incidents d'alimentation ;
    \item les déformations de cartouches ;
    \item les défauts de positionnement.
\end{itemize}

\section{Extraction}

Après le tir, l'étui doit être retiré de la chambre.

Cette fonction est assurée par l'extracteur.

L'extraction peut devenir difficile lorsque :

\begin{itemize}
    \item les pressions sont élevées ;
    \item la chambre est encrassée ;
    \item l'étui est endommagé.
\end{itemize}

\section{Éjection}

Une fois extrait, l'étui doit être expulsé de l'arme.

Cette opération est généralement assurée par un éjecteur.

L'éjection doit être suffisamment énergique pour éviter toute
interférence avec le cycle suivant.

\section{Réarmement}

Le réarmement prépare l'arme pour le tir suivant.

Selon le système considéré, cette opération peut être :

\begin{itemize}
    \item manuelle ;
    \item semi-automatique ;
    \item automatique.
\end{itemize}

\section{Modes de fonctionnement}

\subsection{Répétition manuelle}

Chaque cycle est initié par le tireur.

Exemples :

\begin{itemize}
    \item fusils à verrou ;
    \item armes à levier ;
    \item armes à pompe.
\end{itemize}

\subsection{Semi-automatique}

L'énergie du tir assure le cycle mécanique.

Le tireur doit toutefois actionner la détente pour chaque coup.

\subsection{Automatique}

Le cycle complet est répété tant que :

\begin{itemize}
    \item la détente est maintenue ;
    \item des munitions sont disponibles.
\end{itemize}

\section{Sources d'incidents}

Les incidents d'alimentation peuvent avoir plusieurs origines :

\begin{itemize}
    \item chargeur défectueux ;
    \item bande endommagée ;
    \item ressort fatigué ;
    \item encrassement ;
    \item cartouche déformée ;
    \item défaut mécanique.
\end{itemize}

Ces incidents figurent parmi les pannes les plus courantes des armes
modernes.

\section{Influence sur les performances}

Le système d'alimentation influence notamment :

\begin{itemize}
    \item la cadence pratique ;
    \item la cadence théorique ;
    \item la fiabilité ;
    \item la disponibilité opérationnelle.
\end{itemize}

Pour les armes automatiques, ces paramètres sont souvent aussi importants
que les performances purement balistiques.

\section{Application aux simulations}

Dans de nombreuses simulations, l'alimentation est modélisée de manière
simplifiée.

Les paramètres les plus courants sont :

\begin{itemize}
    \item capacité du chargeur ;
    \item cadence de tir ;
    \item temps de rechargement ;
    \item probabilité de panne.
\end{itemize}

Les simulateurs les plus avancés représentent explicitement le cycle de
fonctionnement et certains incidents.

\begin{simulation}

Dans les simulateurs d'entraînement, les incidents d'alimentation peuvent
être volontairement introduits afin de former les utilisateurs aux
procédures de dépannage.

\end{simulation}

\section{Vers les chapitres suivants}

Après avoir étudié la structure mécanique de l'arme, nous pourrons nous
intéresser aux interfaces entre l'arme et le tireur :

\begin{itemize}
    \item organes de visée ;
    \item viseurs optiques ;
    \item systèmes électroniques ;
    \item conduites de tir.
\end{itemize}

\section{À retenir}

\begin{aretenir}

\begin{itemize}
    \item Le système d'alimentation approvisionne l'arme en munitions.
    \item Le cycle de fonctionnement comprend plusieurs étapes successives.
    \item L'extraction et l'éjection sont essentielles à la continuité du
    tir.
    \item La fiabilité dépend fortement de la qualité du système
    d'alimentation.
    \item Les incidents d'alimentation figurent parmi les pannes les plus
    fréquentes.
\end{itemize}

\end{aretenir}
